% MeshCore for Bitcoiners --- a Bitcoin Austria talk.
% Build from the repo ROOT:  make           (or: latexmk -xelatex meshcore-for-bitcoiners.tex)
% Requires the `theme/` submodule (Bitcoin Austria beamer style). If it is
% missing, run:  git submodule update --init --recursive
%
% NOTE: this is a working draft built from handoff-20260626.md; develop further.
% Spoken cues, staging, and fact-check provenance live in speakers-notes.md ---
% update it in the same commit when a slide's facts or staging change.
%
% Aspect ratio is parameterised so the deck can be built at 16:9 (default) or
% 4:3 without editing this file: the makefile passes `-usepretex='\def\baAspect{43}'`.
% \providecommand keeps 169 as the default for a plain `latexmk`/`make` run.
\providecommand{\baAspect}{169}
\documentclass[aspectratio=\baAspect,t,9pt]{beamer}

\usepackage{theme/bitcoin-austria}

% Talk-specific helper: a Bitcoin vs MeshCore comparison with a "where the
% analogy breaks" callout (wraps the theme's generic \comparisonslide).
% Pass an empty 4th argument {} to omit the callout.
\newcommand{\bridgeslide}[4]{%
  \comparisonslide{#1}{Bitcoin}{#2}{MeshCore}{#3}{Where it breaks}{#4}}

% Talk-specific: app / website screenshots, to break up the theory with a look
% at the real thing. Phone screenshots are tall (portrait), so each is shown as
% a "card" --- the image with a short caption beneath. Cards are narrow, so two
% (or more) sit side by side on one slide; add more as they arrive.
%   \screenshotcard[width]{image}{caption}
\newcommand{\screenshotcard}[3][0.30\paperwidth]{%
  \begin{minipage}[t]{#1}%
    \centering
    \includegraphics[height=0.50\paperheight,keepaspectratio]{#2}\par
    \vskip0.6em
    {\footnotesize #3\par}%
  \end{minipage}%
}

% Theme extras used in this deck: \cornerimage (floats a small image into a
% slide corner, clear of the headline band; call it right before the slide) and
% \narrowonly{<code>} (runs <code> only on narrow/4:3 builds --- the escape
% hatch for a slide still too dense after the theme's auto-scaling). Footnote
% styling (the TOFU / PSK glosses) also comes from the theme.

\title{MeshCore for Bitcoiners}
\subtitle{Permissionless communication --- from money to radio}
\author{Harald Schilly --- Bitcoin Austria}
\date{\today}

\begin{document}

\begin{frame}[plain,noframenumbering]
  \titlepage
\end{frame}

% ============================================================ what is MeshCore?
% Bitcoiners don't know MeshCore --- introduce it before drawing the parallels.
\fillerslide[What Bitcoin did for money, MeshCore does for messaging!]{What is MeshCore?}

% Hook slide: open with the one-line mental model + why a Bitcoiner should care,
% BEFORE the mechanics. Spoken opener (the stakes question): speakers-notes.md.
\begin{frame}{The one-line version}
  \begin{itemize}
    \item What Bitcoin did for money, MeshCore does for messaging:
      \textbf{a radio mesh you can join without asking}
    \item Generate a \textbf{keypair} and you're on --- no carrier, no number,
      no SIM, nobody to ask for permission
    \item Cheap radios relay each other in a \textbf{mesh}; coverage grows because
      people join, not because a company builds towers
  \end{itemize}
  \vfill
  \begin{alertblock}{Why a Bitcoiner should care}
    Spectrum is licensed, telcos are KYC gatekeepers, and the network itself is
    the chokepoint. MeshCore routes around all three ---
    \textbf{censorship-resistant, self-sovereign comms that still work when the
    grid is down}.
  \end{alertblock}
\end{frame}

\begin{frame}{Messaging without the network}
  \begin{itemize}
    \item Off-grid text messaging over long-range \textbf{LoRa} radio ---
      no internet, no cell tower, no SIM
    \item \textbf{Open-source firmware} running on cheap LoRa boards
    \item Uses \textbf{licence-free} ISM radio bands --- worldwide, region by
      region (EU 868\,MHz; US/AU 915\,MHz; \dots)
  \end{itemize}
  \vfill
  \begin{alertblock}{The EU angle}
    On these bands \emph{anyone} may transmit --- no licence, no registration, no
    fee. You only respect the power and duty-cycle limits.
  \end{alertblock}
\end{frame}

% Show the real thing early: a packet that actually routed across Austria, so
% the audience sees it works before we get into the how.
\fullbleedslide[A packet routed across Austria \;\textbullet\; logger-at.meshcore.observer]{observer-package-across-austria-20260626}

% The Bitcoiner's reflex after seeing it work: "so what's the catch?" Name the
% range/bandwidth trade-off now; the throughput bridge later pays it off.
\begin{frame}{The catch --- kilometres of range, bytes of bandwidth}
  \begin{itemize}
    \item \textbf{Long range, tiny pipe:} one \emph{hop} reaches \emph{tens} of
      kilometres line-of-sight (the packet above crossed Austria by chaining
      many) --- but at roughly the speed of a 1980s modem. One LoRa frame is at
      most \textbf{255 bytes}
    \item \textbf{Store-and-forward, not a phone call:} messages hop repeater to
      repeater with seconds of latency; you text, you don't stream
    \item \textbf{Physics, not policy:} the range buys reach by spending
      bandwidth and time --- and it sips power, so a node runs for days off a
      small battery, or indefinitely on a small solar panel
  \end{itemize}
  \vfill
  \begin{alertblock}{The right mental model}
    Not WhatsApp-over-radio --- think \textbf{a permissionless telegraph}: short
    text, long reach, no carrier. The narrow pipe is the price of needing nothing
    but a cheap radio and line of sight.
  \end{alertblock}
\end{frame}

% The 2026 fork is on-theme for a Bitcoin crowd: MIT firmware, a brand-capture
% attempt, the community routing around it. We stick to the core team's .io
% domain throughout. Facts verified --- details, sources, and the spoken
% punchline lead-in: speakers-notes.md ("Where it comes from").
\begin{frame}{Where it comes from}
  \begin{itemize}
    \item Grew out of the \textbf{Meshtastic} world, then rebuilt from scratch
      --- a separate network, \emph{not} interoperable
    \item Created by \textbf{Scott Powell} in \textbf{late 2024};
      web/mobile clients by Liam Cottle
    \item \textbf{2026 split:} promoter Andy Kirby quietly filed for the
      \textbf{MeshCore trademark} and pushed undisclosed AI-rewritten apps ---
      the core team broke with him
    \item \textbf{MIT-licensed}, so the grab couldn't capture the protocol: the
      team kept shipping at \textbf{meshcore.io} (Kirby kept the old
      \texttt{.co.uk})
    \item \textbf{Still open, as we speak:} the community \textbf{crowdfunded
      the legal defence} in July, the opposition is filed --- the mark sits in
      \textbf{opposed} status, awaiting a ruling
  \end{itemize}
  \vfill
  \begin{alertblock}{A familiar story}
    Permissionless tech, a \textbf{brand-capture attempt} --- the
    \texttt{.co.uk} holder trademarking the name, much as \textbf{bitcoin.com}
    marketed Bitcoin Cash as ``Bitcoin'' --- and the community just
    \textbf{routes around it}: the code stayed free, so the fight is only over
    \emph{a name}. Bitcoiners have seen this movie.
  \end{alertblock}
\end{frame}

\begin{frame}{The hardware --- tens of euros, not hundreds}
  \begin{itemize}
    \item \textbf{Companion} (your device: screen + Bluetooth) --- SenseCAP
      T1000-E, credit-card size with GPS, $\sim$40\,EUR
    \item \textbf{Portable with a display} --- Heltec LoRa\,32 V4 (OLED,
      $\sim$40\,EUR), Seeed Wio Tracker L1 Pro (screen, GPS, case)
    \item \textbf{Repeater} --- SenseCAP Solar Node P1-Pro: 5\,W solar +
      battery, runs 24/7 off-grid, $\sim$100\,EUR $+$ $\sim$30 antenna
    \item \textbf{Build your own} --- Heltec / RAK WisBlock boards
  \end{itemize}
  \vfill
  % A small device row, in bullet order: companion / portable / repeater.
  % Images are background-removed, tightly trimmed PNGs; height-capped so they
  % align and the row stays within \textwidth on both 16:9 and 4:3.
  \centering
  \begin{minipage}[b]{0.27\paperwidth}\centering
    \includegraphics[height=0.28\paperheight,width=\linewidth,keepaspectratio]{hw-t1000e}%
  \end{minipage}\hfill
  \begin{minipage}[b]{0.27\paperwidth}\centering
    \includegraphics[height=0.28\paperheight,width=\linewidth,keepaspectratio]{hw-heltec-v4}%
  \end{minipage}\hfill
  \begin{minipage}[b]{0.27\paperwidth}\centering
    \includegraphics[height=0.28\paperheight,width=\linewidth,keepaspectratio]{hw-sensecap-p1pro}%
  \end{minipage}
\end{frame}

% Practical example, right after "the hardware is cheap": the four steps to
% actually get on the mesh that night --- buy, pair, generate, join.
\begin{frame}{Your first five minutes on the mesh}
  \begin{enumerate}
    \item \textbf{Get a companion device} --- e.g.\ a SenseCAP \textbf{T1000-E}
      ($\sim$40\,EUR)
    \item \textbf{Flash it} --- connect over USB and flash from the browser at
      \href{https://flasher.meshcore.io/}{flasher.meshcore.io}
    \item \textbf{Install the app} --- Liam Cottle's MeshCore client
      (iOS / Android / web) and pair the board over \textbf{Bluetooth}
    \item \textbf{It generates your keypair on first run} --- that pair \emph{is}
      your identity. Back it up; nobody else can recover it
    \item \textbf{Send your first message} --- to the public \textbf{\#test}
      channel: no account, no SIM, no permission
  \end{enumerate}
  \vfill
  \begin{alertblock}{That's the whole onboarding}
    Buy, flash, pair, generate, send --- the permissionless claim is literal:
    nobody issued you anything, and nothing on the network knew you were coming.
  \end{alertblock}
\end{frame}

% MeshCore (red, the subject) vs Meshtastic --- generic comparison, not \bridgeslide
\comparisonslide{MeshCore vs Meshtastic}%
  {MeshCore}{%
    \begin{itemize}
      \item \textbf{Role-based}: companion (handset), repeater (always-on
        relay), room server (store-and-forward)
      \item \textbf{Source-routed} along discovered paths
      \item Quieter, \textbf{more scalable mesh}
    \end{itemize}}%
  {Meshtastic}{%
    \begin{itemize}
      \item Easy to join, \textbf{phone-first}
      \item \textbf{Flood-based} --- every node rebroadcasts
      \item Gets \textbf{noisy} as it grows
    \end{itemize}}%
  {Same cheap hardware}{%
    Same boards, same chips --- a different firmware and routing philosophy.}

% Break up the theory with a look at the real app: watch it, chat, plan a link.
% Staging (switch to the live observer here if connectivity allows): speakers-notes.md.
\begin{frame}{Live on the mesh}
  \vfill
  \centering
  \screenshotcard[0.27\paperwidth]{meshcore-observer-live-map-vienna}{%
    The observer's live map: 828 nodes, packets arcing across Vienna in
    real time.}%
  \hspace{1em}%
  \screenshotcard[0.27\paperwidth]{meshcore-app-test-channel}{%
    The public \textbf{\#test} channel: signed adverts and \emph{ping}\,/\,%
    \emph{pong} probes scrolling by.}%
  \hspace{1em}%
  \screenshotcard[0.27\paperwidth]{meshcore-app-line-of-sight-vienna}{%
    Line-of-sight planner: terrain and Fresnel zone between two repeaters
    (6.7\,km in Vienna).}%
  \vfill
\end{frame}

% ------------------------------------------------------------------- routing
% The mechanism behind "source-routed", unravelled: flood to discover, then
% follow the remembered path. Two slides: the DM lifecycle, then what always
% floods. Source: kiekr.app/meshcore-explained; details in speakers-notes.md
% ("Routing").
\begin{frame}{Routing --- flood once, then follow the path}
  \begin{enumerate}
    \item \textbf{First contact floods} --- no path exists yet, so every
      repeater relays the message once, stamping its own ID into the packet
      as it passes
    \item \textbf{The ACK carries the path back} --- the delivery confirmation
      returns with the accumulated route, and the sender remembers it
    \item \textbf{Later messages follow the path} --- source-routed repeater by
      repeater, no flooding (each direction keeps its own path)
    \item \textbf{Stale path? Flood again} --- a failed route falls back to
      flooding and the path is re-learned automatically
  \end{enumerate}
  \vfill
  \begin{alertblock}{Where the routing table lives: nowhere}
    Repeaters relay and forget --- the \textbf{sender holds the path}, per
    contact and per direction. Discovery is \emph{reactive}: pay for one
    flood, then ride the discovered path for free.
  \end{alertblock}
\end{frame}

\begin{frame}{What always floods}
  \begin{itemize}
    \item \textbf{Channel messages always flood} --- a channel is just a
      shared key with no member list, so there is no single recipient to
      route to
    \item \textbf{And nobody ACKs them} --- no delivery confirmation exists,
      whoever is offline simply misses the message
    \item \textbf{Adverts flood too} --- the signed ``business card'' (name,
      public key, position) goes \emph{zero-hop} to direct neighbours, or
      region-wide via every repeater
    \item \textbf{Flooding is the expensive mode} --- every repeater in reach
      spends its scarce airtime on the packet, relevant to its users or not
  \end{itemize}
  \vfill
  \begin{alertblock}{This is why regions exist}
    A flood stops only where repeaters stop relaying. In a moment: how
    MeshCore \textbf{fences the flood by agreement} --- the region code.
  \end{alertblock}
\end{frame}

% The flip side of source routing: a real message and the path it travelled.
% Content is top-aligned (not vfill-centred) so the captions sit directly under
% the images, not pushed down to the slide edge.
\begin{frame}{How a message reaches me}
  \vspace{0.8em}
  \centering
  \begin{minipage}[t]{0.44\paperwidth}%
    \centering
    \includegraphics[width=0.44\paperwidth,keepaspectratio]{meshcore-app-coffee-message}\par
    \vskip0.7em
    {\footnotesize A single message on the public \textbf{\#coffee} channel,
     opened to show its metadata: path-hash size, SNR, and the hop count it
     picked up crossing the mesh.}%
  \end{minipage}%
  \hspace{2.5em}%
  \begin{minipage}[t]{0.30\paperwidth}%
    \centering
    \includegraphics[height=0.46\paperheight,keepaspectratio]{meshcore-app-path-7hops}\par
    \vskip0.7em
    {\footnotesize \dots\ and a delivery path traced hop by hop: each
     \emph{known repeater} that relayed the message, from sender to me.}%
  \end{minipage}%
  \vfill
\end{frame}

% --------------------------------------------------------------- region codes
% The region "transport code" is a 2-byte HMAC of the region key over the
% (encrypted) payload, recomputed per hop --- geofencing the flood. Construction
% verified against the MeshCore source; details & sources: speakers-notes.md.
\fillerslide[Borders agreed by name, enforced per hop]{The region transport code}

% Plain-language setup before the bridge: the audience doesn't know MeshCore, so
% explain what a region code is and why it exists before drawing the SIGHASH bridge.
\begin{frame}{What is a region code, and why?}
  \begin{columns}[T]
    \begin{column}{0.66\textwidth}
      \begin{itemize}
        \item A \textbf{region} is just an agreed name --- e.g.\ \textcolor{baRed}{\texttt{at}}
          for all of Austria, or \textcolor{baRed}{\texttt{at-west}} / \textcolor{baRed}{\texttt{at-ost}}
          / \textcolor{baRed}{\texttt{at-sued}} for its three macro-regions (right)
        \item Every packet carries a 2-byte \textbf{transport code} --- a hash
          over region name and message body: repeaters relay only on a match
          with their own regions
        \item The mesh is thus \textbf{geofenced by agreement} --- no registry,
          no coordinates, no permission: just a name the locals share
        \item Usage rule: \textbf{pick the smallest region} that covers your
          recipients --- village chatter posted \texttt{at}-wide burns the
          whole country's airtime
      \end{itemize}
    \end{column}
    \begin{column}{0.32\textwidth}
      \centering
      \includegraphics[width=\linewidth]{austria-regions}\\[0.3em]
      {\scriptsize Austria's three MeshCore macro-regions ---
       \texttt{at-west} (west) / \texttt{at-ost} (east) / \texttt{at-sued} (south);
       cities nest below, e.g.\ \texttt{at-w} = Vienna\\
       region scheme: \href{https://meshcore-austria.at}{meshcore-austria.at}}
    \end{column}
  \end{columns}
  \vfill
  \begin{alertblock}{The neat part}
    The code is computed over the \emph{encrypted} payload, so repeaters filter by
    region \textbf{without decrypting} --- and because it \textbf{commits to the
    packet's bytes}, you can't lift it onto a different packet (a Bitcoin parallel
    we'll sharpen later).
  \end{alertblock}
\end{frame}

% One statement of the precise claim + the primitives preview + the honesty
% contract (merged from three earlier restatements of the same thesis).
\begin{frame}{The precise claim --- and a promise}
  We opened with the one-liner; here is the claim the rest of the talk defends ---
  the same Bitcoin primitives, in radio:
  \vskip0.6em
  \begin{itemize}
    \item A self-generated \textbf{keypair replaces institutional identity} ---
      Bitcoin fires the bank; MeshCore fires the telco \emph{and} the SIM
      registry behind it
    \item \textbf{Signatures and encryption} do the work --- adverts signed,
      DMs encrypted
    \item \textbf{A throughput cap you didn't choose} forces a second layer
    \item \textbf{Permissionless by construction} --- routing not flooding, and
      you join by generating a key, not by asking
  \end{itemize}
  \vfill
  \begin{alertblock}{The line the talk hangs on}
    A permissionless, self-issued identity broadcasting in the ether is
    \textbf{arguably even more subversive than permissionless money}. Every time
    the parallel \emph{breaks}, I'll say so out loud --- the breaks are the
    interesting part.
  \end{alertblock}
\end{frame}

% ---------------------------------------------------------------- identity
\fillerslide[The strongest parallel]{Identity is a keypair}

\bridgeslide{Identity is a keypair}{%
  \begin{itemize}
    \item \textbf{Self-generated keypair}, no issuer
    \item Address = \textbf{hash of the public key}
    \item \textbf{``Not your keys, not your coins''}
    \item Verify the address \textbf{out-of-band}
  \end{itemize}}{%
  \begin{itemize}
    \item \textbf{Self-generated keypair}, on-device
    \item Routing target = \textbf{public-key prefix}
    \item \textbf{``Not your keys, not your identity''}
    \item Verify keys via \textbf{QR / word list} (TOFU\footnote{Trust On First Use --- check a key once out-of-band, then pin it.})
  \end{itemize}}{%
  Two different tricks for the short ID: Bitcoin \textbf{hashes} the pubkey (hidden
  until you spend --- Taproot skips the hash); MeshCore just \textbf{truncates} it
  to a prefix --- a routing shortcut, so it can collide (1 byte originally, 1--3
  today). And the identity is deliberately
  \emph{sticky} (no HD rotation, no seed phrase) --- so \textbf{back it up}: lose
  the key and you lose the \emph{handle}.}

% ---------------------------------------------------------------- signing
\fillerslide[The elegant twist]{The signing model}

\bridgeslide{Opposite jobs for the same keypair}{%
  \begin{itemize}
    \item Content is \textbf{public} (on-chain)
    \item The key protects \textbf{authorization}
    \item No valid signature $\to$ the transaction \textbf{does not exist}
  \end{itemize}}{%
  \begin{itemize}
    \item Identity is \textbf{public} (the advert --- a node's signed
      ``here is my key, here I am'')
    \item The key protects \textbf{content}
    \item Adverts are \textbf{signed} (Ed25519)
    \item DMs are \textbf{encrypted} (X25519)
  \end{itemize}}{%
  But channels/rooms share one PSK\footnote{Pre-Shared Key --- one secret everyone
  in the channel holds.}, so a group message proves only that \emph{someone} in the
  channel sent it --- \textbf{not who}. Anyone with the PSK can read \emph{and}
  forge the sender name. DMs (X25519) and the advert (Ed25519) authenticate;
  \textbf{channel posts don't}.}

% Region code <-> SIGHASH (moved here from the intro region explanation): the
% 2-byte transport code commits to the (encrypted) payload, like a signature
% commits to its own tx. Audience already met region codes in the intro.
\bridgeslide{The region rides inside the message}{%
  \begin{itemize}
    \item A signature \textbf{commits to the transaction} it signs --- SIGHASH
      flags even pick \emph{which} parts it covers
    \item Lift it onto a different tx and it is invalid --- authorization is
      \textbf{bound to the content}
  \end{itemize}}{%
  \begin{itemize}
    \item The 2-byte region code is an \textbf{HMAC\footnote{Hash-based Message
      Authentication Code --- a keyed hash of the message.} committing to the
      payload}; repeaters recompute it per hop to decide whether to relay
    \item The region key is just \texttt{SHA256(name)[:16]} --- agree on
      ``\texttt{at-w}'' (Vienna), no key exchange, no registry
  \end{itemize}}{%
  No flags, no signer, no authorship: symmetric and \textbf{forgeable} (anyone who
  knows the region name can recompute it) and only 2 bytes wide --- a flood-scoping
  filter that saves airtime. It binds \emph{scope}, not \emph{authorship}.}

% ---------------------------------------------------------------- throughput
\fillerslide[A chosen constraint, not a weakness]{Throughput}

% Numbers verified (2026-07-02), both sides: EU band-P duty cycle (10%, 500 mW
% ERP), the EU/UK Narrow preset 869.618 MHz / SF8 / BW62.5 / CR4/8, the LoRa
% time-on-air figures, and the Bitcoin tx/block rate. Derivations, sources, Q&A
% nuances, and the "six seconds" reveal: speakers-notes.md ("Throughput").
% DON'T revert the frequency to 869.525 --- that's the old wide preset (the
% LoRaWAN RX2 frequency), not the Narrow preset.
\bridgeslide{The throttle is imposed, not a failure}{%
  \begin{itemize}
    \item Block space is a \textbf{consensus} cap (4M weight units): one block
      every $\sim$10\,min
    \item $\sim$3{,}000--4{,}000 tx/block $\Rightarrow$ $\sim$5--7\,tx/s
      $\approx$ \textbf{300--400 tx/min}, network-wide
  \end{itemize}}{%
  \begin{itemize}
    \item Duty cycle is a \textbf{regulatory} cap: 10\% airtime on the
      869.4--869.65\,MHz sub-band (EU default: 869.618\,MHz, SF\,8, BW\,62.5\,kHz, CR\,4/8)
    \item One msg $\approx$ 0.25--1.6\,s on air (short text $\to$ full 184\,B);
      10\% $=$ 6\,s airtime/min $\Rightarrow$ \textbf{$\sim$4--24 msgs/min per node}
  \end{itemize}}{%
  Bitcoin's rate is one \emph{global} ledger; MeshCore's is \emph{per node} on a
  \textbf{shared} channel --- many nodes split the same airtime, so the real
  ceiling is lower. Either way the throttle is imposed, not ``they couldn't do
  more'' --- and a deliberate cap is exactly what forces a second layer to exist.}

% ---------------------------------------------------------------- lightning
\fillerslide[Two distinct analogies --- keep them separate]{Lightning Network}

\bridgeslide{Angle A --- route along a discovered path}{%
  \begin{itemize}
    \item \textbf{Route through specific channels} --- you pick the path
    \item Don't broadcast the payment \textbf{to the whole network}
  \end{itemize}}{%
  \begin{itemize}
    \item An advert floods once; \textbf{path info accumulates}
    \item Later DMs are \textbf{source-routed along that path}
  \end{itemize}}{%
  \begin{itemize}
    \item \textbf{Discovery differs:} Lightning gossips the whole \emph{public}
      graph (proactive); MeshCore floods one advert to learn \emph{one} path
      (reactive)
    \item Lightning \textbf{onion-routes} (hop privacy); MeshCore path bytes ride
      in the clear
    \item Lightning carries \textbf{atomic value} (HTLCs); MeshCore carries data,
      no settlement
  \end{itemize}}

\begin{frame}{Angle B --- bootstrap, then jump layers}
  \begin{itemize}
    \item Move the bulk transfer to a \textbf{fat pipe} (internet, Wi-Fi,
      sneakernet) --- off-mesh
    \item The mesh carries only the rare, \textbf{expensive anchoring event}, never the payload
    \item Key and rendezvous get negotiated over \textbf{encrypted DMs} --- on the mesh itself
  \end{itemize}
  \vfill
  \begin{alertblock}{Where it breaks}
    Lightning stays cryptographically chained to L1 (you can unilaterally
    close). The mesh\,→\,fat-pipe handoff is a \textbf{one-time bootstrap}, not
    a trust-anchored channel --- closer to WebRTC signaling or a BitTorrent
    tracker.
  \end{alertblock}
\end{frame}

% ---------------------------------------------------------------- facts
\fillerslide[Side by side --- what's under the hood]{Crypto facts}

\begin{frame}{The crypto, side by side}
  \centering
  \comparisontable{Bitcoin}{MeshCore}{%
    Curve(s)      & secp256k1               & Ed25519 + X25519 \\
    Identity      & address = hash(pubkey)  & pubkey prefix (1--3\,B) \\
    Signatures    & every transaction       & adverts only \\
    Encryption    & none (all public)       & X25519 (DMs), AES-128 (groups) \\
    The key's job & authorize spending      & protect content / identity \\
    Address reuse & rotated (HD wallets)    & sticky, stay discoverable \\
  }

  \vspace{0.8em}
  {\footnotesize Routing prefix: 1 byte originally (collisions); 1--3 bytes
  since firmware 1.14+.}
\end{frame}

% ---------------------------------------------------------------- channel hash
% A Bitcoiner spots SHA256-twice and reaches for "mining" --- same primitive,
% different game. Construction verified against the MeshCore source (AES-128-ECB
% + 2-byte MAC, NOT AES-256/CTR); details & sources: speakers-notes.md.
\fillerslide[Same primitive, a different game]{Double SHA256 --- but not mining}

\begin{frame}{Joining \texttt{\#bitcoin-at}: SHA256, twice, cropped}
  \texttt{SHA256(SHA256(\#name)[:16])} --- \textbf{hash} the channel
  \texttt{\#name}, \textbf{crop} to 16 bytes (the AES-128 key), \textbf{hash
  again}, take the first byte as a tiny channel id. That \textbf{public id rides
  in the clear} --- it narrows a receiver to a couple of candidate keys, and the
  MAC\footnote{Message Authentication Code --- a keyed tag; only the matching key
  produces a valid one.} confirms the match. Two SHA256 rounds \emph{look} like
  Bitcoin --- but this is not mining.
  \vfill
  \centering\small
  \setlength{\tabcolsep}{8pt}\renewcommand{\arraystretch}{1.3}%
  \begin{tabular}{@{}lccc@{}}
    \toprule
      & \textcolor{baRed}{\textbf{MeshCore}} & \multicolumn{2}{c}{\textbf{Bitcoin}} \\
    \cmidrule(lr){3-4}
      & channel hash & txid / block id & \textbf{mining} \\
    \midrule
      Construction      & SHA256 + crop + SHA256 & SHA256d (full) & SHA256d (full) \\
      Computed          & once, fixed            & once, fixed    & \textbf{nonce search} \\
      Difficulty target & no                     & no             & \textbf{yes} \\
      Purpose           & id / commitment        & id / commitment & proof-of-work \\
    \bottomrule
  \end{tabular}
  \vfill
  \begin{alertblock}{The parallel that holds}
    No nonce search, no target, no contest --- so not mining. The real Bitcoin
    echo is \textbf{a short one-way commitment that trades length for
    collisions} --- like Bitcoin's 160-bit HASH160 addresses or MeshCore's
    1--3-byte routing prefixes.
  \end{alertblock}
\end{frame}

% --------------------------------------------------------------- ASIC / silicon
% Follows "not mining": a Bitcoiner asks whether LoRa has ASICs too. It does ---
% the SX1262 runs the LoRa PHY in silicon, but for power, not a PoW contest.
% Facts & sources: speakers-notes.md ("ASICs --- and the SX1262").
\fillerslide[Same silicon idea, opposite reason]{ASICs --- and the LoRa radio chip}

% Float a small photo of the SX1262 high in the top-right corner of the
% bridgeslide --- small enough to sit up by the title/page number, not over the
% MeshCore column. (Bare name; resolves via the deck's pix/ graphics path.)
\cornerimage[width=0.10\paperwidth,topclear=20pt]{sx1262}
\bridgeslide{Purpose-built silicon, for opposite reasons}{%
  \begin{itemize}
    \item Mining \textbf{ASICs} race to compute SHA256 hashes
    \item Efficiency wins \textbf{reward} --- an arms race; more hashpower is
      always better, and difficulty just pulls up to match
  \end{itemize}}{%
  \begin{itemize}
    \item The \textbf{SX1262} radio chip runs the whole LoRa modem in hardware,
      not in software
    \item Dedicated silicon for \textbf{power}: microamps asleep, so a small
      battery runs a node for days --- an SDR \emph{can} decode LoRa in
      software, just at orders of magnitude more power
  \end{itemize}}{%
  Both are purpose-built chips --- but Bitcoin's ASIC wins a \emph{contest} (more
  efficiency $\to$ more reward), while the SX1262 just \emph{sips power} at a
  fixed job. No race to win: the ceiling is physics and the duty cycle, not rival
  hashpower.}

% ---------------------------------------------------------- brainwallet rooms
% A #hashtag room's AES-128 key is SHA256 of its name --- no salt, no
% stretching: a brainwallet. The ~90 s crack figure and the WebGPU cracker are
% verified; details & sources: speakers-notes.md ("Hashtag rooms").
\fillerslide[The name is the key --- a familiar footgun]{Hashtag rooms are brainwallets}

\begin{frame}{A hashtag room is a brainwallet}
  We just saw the room key \emph{is} \texttt{SHA256} of the name --- a Bitcoiner
  has seen exactly this go wrong.
  \vskip0.3em
  \centering
  \comparisontable{Bitcoin brainwallet}{MeshCore \#room}{%
    Key from           & \texttt{SHA256(passphrase)} & \texttt{SHA256(\#name)[:16]} \\
    Input              & low-entropy phrase          & short word (\texttt{\#vienna}) \\
    Salt / stretching  & none                        & none \\
    Public oracle      & the blockchain              & the broadcast packet \\
    Failure            & swept in seconds            & cracked in $\sim$90\,s ($<$7 chars) \\
  }
  \vfill
  \begin{alertblock}{The deep point --- a free verification oracle}
    Both hand the attacker a public way to \emph{check} an offline guess: the chain
    says ``this address holds coins''; the broadcast says ``this key decodes a
    message.'' Offline guessing \emph{plus} public verification is what makes both
    catastrophic --- \textbf{not} an AES break: the entropy is the \emph{name}.
  \end{alertblock}
\end{frame}

\begin{frame}{\dots\ but public \emph{by design}}
  \begin{itemize}
    \item \textbf{Lower stakes:} a brainwallet loses \emph{value}, irreversibly;
      a hashtag room loses \emph{confidentiality} (eavesdropping), not coins.
    \item \textbf{And it flatters MeshCore:} hashtag rooms are labelled
      \textbf{world-readable by design} --- ``no key to remember'' is the feature,
      not a hidden footgun. MeshCore \emph{tells} you the room is public.
    \item \textbf{Know which layer is public:} private traffic gets real crypto
      --- X25519\,/\,Ed25519 DMs, or private channels with long random hex keys.
  \end{itemize}
  \vfill
  \begin{alertblock}{The line}
    You wouldn't put your coins in \texttt{SHA256(``correct horse'')}. Don't put
    your secrets in \texttt{\#vienna}.
  \end{alertblock}
\end{frame}

% ============================ what Bitcoin's consensus buys, and MeshCore skips
% A four-slide section flowing from one insight: Bitcoin defends a scarce asset
% (coins), so it pays for heavy consensus machinery (PoW/Sybil resistance, one
% global order, permanent state); MeshCore defends nothing scarce -- it relays
% fleeting messages -- so it needs none of it. Each slide is one face of that.
% (Per-topic divider slides folded into the bridge titles.)
\fillerslide[What is all that consensus machinery actually for?]{What Bitcoin's consensus buys --- and MeshCore skips}

% #1 Sybil resistance.
\bridgeslide{Sybil resistance --- what proof-of-work buys}{%
  \begin{itemize}
    \item PoW makes identity / participation deliberately \textbf{costly}
    \item That cost \textbf{is the security}: Sybil resistance, global ordering,
      permanent state
    \item You spend energy to earn a say
  \end{itemize}}{%
  \begin{itemize}
    \item Identity is \textbf{free} --- an Ed25519 keypair, generated offline, in
      unlimited number
    \item \textbf{No cost to participate} --- no proof-of-work
    \item Anyone can flood adverts
  \end{itemize}}{%
  \textbf{No Sybil resistance --- and none needed:} nothing to vote on, no scarce
  asset to seize. Abuse is only \emph{throttled}, never stopped --- \textbf{physical}
  limits (airtime / duty cycle, RF range) and \textbf{social} ones (repeater ACLs,
  admin passwords). The worst case is wasted airtime, not theft.}

% #2 Incentives.
\bridgeslide{Incentives --- paid security vs.\ a volunteer commons}{%
  \begin{itemize}
    \item Miners are \textbf{paid} --- block subsidy + fees
    \item Scarce block space $\to$ a \textbf{fee market} (a price)
    \item Incentives solve the free-rider problem
  \end{itemize}}{%
  \begin{itemize}
    \item Repeater operators are \textbf{unpaid} --- pure goodwill
    \item Scarce airtime $\to$ \textbf{social etiquette}: don't spam Public,
      use \#test, keep advert intervals conservative
    \item No price, no reward
  \end{itemize}}{%
  Bitcoin engineers incentives so the network sustains itself without trust;
  MeshCore \textbf{leans on volunteers and norms --- a commons with no fee
  market}. It works at today's scale, but nothing \emph{economic} defends it
  against the tragedy of the commons.}

% #3 State.
\bridgeslide{State --- consensus state vs.\ the moment}{%
  \begin{itemize}
    \item A \textbf{permanent, globally agreed ledger}
    \item Confirmed once $\to$ \textbf{there forever}, verifiable by anyone, later
  \end{itemize}}{%
  \begin{itemize}
    \item Channel messages exist \textbf{only as they are transmitted}
    \item Offline when it's sent $\to$ you miss it, silently, with no indication
    \item Room servers add \emph{optional} short-term history
  \end{itemize}}{%
  Bitcoin's whole purpose is durable shared state; MeshCore channels deliberately
  keep none. Different goals --- \textbf{you can't ``re-sync'' a missed channel
  message} the way a node replays the chain from genesis.}

% #4 Time. The forward-in-time rule is scoped to ADVERTS (per-pubkey
% last_advert_timestamp), NOT to every packet --- keep the slide's wording.
% Verification & sources: speakers-notes.md ("Time --- advert anti-replay").
\bridgeslide{Time --- two ways to agree on ``after''}{%
  \begin{itemize}
    \item One chain orders \textbf{everything}, globally
    \item That total order is what makes a \textbf{double-spend infeasible}
  \end{itemize}}{%
  \begin{itemize}
    \item A node's \textbf{adverts} must move \textbf{forward in time}: repeaters
      reject any advert not newer than the last one from that key --- a
      per-sender anti-replay rule
    \item Reset your radio with the wrong clock $\to$ your adverts are silently
      dropped and you stop appearing on the mesh
  \end{itemize}}{%
  Bitcoin builds one \emph{global} order (essential against double-spend);
  MeshCore enforces only \emph{per-sender} local order on adverts. There's no
  global ``which came first'' across senders --- with no asset to double-spend,
  it never needs the \emph{global} one.}

% ------------------------------------------------------------------- hard fork
% Firmware 1.14+ widens the routing prefix --- honest caveat: pre-1.14 repeaters
% still relay 1-byte paths (silently dropping wider ones), so old nodes degrade
% rather than drop off; hence "fades off the mesh". Sources: speakers-notes.md.
\fillerslide[Break the wire format, leave the laggards behind]{The multi-byte routing hard fork}

\bridgeslide{A backward-incompatible rule change}{%
  \begin{itemize}
    \item A \textbf{hard fork} loosens the rules, so old nodes reject the new
      blocks
    \item Both sides can stay alive --- \textbf{two live chains}, each kept up by
      its own economic backers
  \end{itemize}}{%
  \begin{itemize}
    \item Firmware 1.14+ \textbf{widens the routing prefix from 1 byte to
      1--3 bytes}
    \item Old repeaters \textbf{silently drop} the wider-prefix packets ---
      they still relay 1-byte paths, but the new traffic passes them by
  \end{itemize}}{%
  No consensus, no ledger, no coins --- nothing keeps a second version alive, so
  old firmware just \textbf{fades off the mesh} as wider prefixes spread. Closer to a
  \textbf{TLS 1.2\,$\to$\,1.3} bump than a chain split: ``hard fork'' names the
  backward incompatibility, not a persistent divergence you could mine on.}

% ===================================================== shortcomings (Part 4)
% Pulls the sharpest breaks into one deliberate beat, split into "deliberate by
% design" vs "real gaps" so the honesty lands without sounding like a teardown.
% Sources: speakers-notes.md ("Shortcomings"). Watch 4:3 fit --- if a slide
% overflows, reach for \narrowonly{\small}.
\fillerslide[The honest reckoning --- what to still distrust]{Shortcomings, as a Bitcoiner sees them}

% Slide 1: the genuine gaps a privacy-minded Bitcoiner cares about most.
\begin{frame}{What still leaks --- even with encryption}
  \begin{itemize}
    \item \textbf{Metadata is in the clear:} adverts broadcast your pubkey and
      rough position, and source-routed paths expose \emph{who relays to whom} ---
      the social graph and location leak even when DM \emph{content} is encrypted
    \item \textbf{No onion routing:} repeaters must read the path prefixes to
      forward, so there is no hop privacy --- unlike Lightning's onion
    \item \textbf{No forward secrecy:} DMs use \emph{static} X25519 to your
      long-term key --- capture that key once and every \emph{recorded past} DM
      decrypts (a Signal-style ratchet is what's missing)
    \item \textbf{Sticky identity widens the blast radius:} the deliberately
      non-rotating key means one compromise exposes \emph{all} of your history
  \end{itemize}
  \vfill
  \begin{alertblock}{The honest line}
    Encryption protects \textbf{what you said}, not \textbf{that you said it, to
    whom, from where} --- and not retroactively if your key ever leaks.
  \end{alertblock}
\end{frame}

% Slide 2: the structural omissions --- mostly deliberate (callbacks), one real.
\begin{frame}{What's left undefended --- and why that's (mostly) on purpose}
  \begin{itemize}
    \item \textbf{No Sybil resistance:} identities are free, so spam is only
      \emph{throttled} by airtime and ACLs --- fine, because there's no scarce
      asset to seize
    \item \textbf{A volunteer commons:} repeaters are unpaid goodwill --- no fee
      market, so nothing \emph{economic} defends the network at scale
    \item \textbf{The default \#public channel ships a well-known key:} everyone
      has it --- convenience, \emph{not} privacy
    \item \textbf{Prefix squatting is cheap:} short routing prefixes are grindable
      --- a nuisance to \emph{discovery}, \emph{not} impersonation (full key
      $+$ TOFU is the backstop)
  \end{itemize}
  \vfill
  \begin{alertblock}{Deliberate vs.\ a real gap}
    Most of these are \textbf{features of defending nothing scarce} --- the worst
    case is wasted airtime, not theft. The two worth fixing before you trust it
    with anything sensitive: \textbf{metadata privacy} and \textbf{forward secrecy}.
  \end{alertblock}
\end{frame}

% Closing call to action: concrete Bitcoin channels the audience can join right
% now. Ties back to the region-transport-code slide (a region is just an agreed
% name: "at" = Austria, "at-w" = Vienna; no region = mesh-wide).
\begin{frame}{Join the Bitcoin channels on MeshCore}
  Same name, same key --- just type the channel name to join.
  \vskip1.2em
  \centering
  \setlength{\tabcolsep}{12pt}\renewcommand{\arraystretch}{1.4}%
  \begin{tabular}{@{}lll@{}}
    \toprule
      \textbf{Channel} & \textbf{Region} & \textbf{Reach} \\
    \midrule
      \texttt{\#bitcoin}      & ---           & whole mesh (no region) \\
      \texttt{\#bitcoin-at}   & \texttt{at}   & Austria \\
      \texttt{\#bitcoin-wien} & \texttt{at-w} & Vienna \\
    \bottomrule
  \end{tabular}
  \vfill
  \begin{alertblock}{Pick your scope}
    Plain \texttt{\#bitcoin} reaches the whole mesh; the regional channels add a
    \textbf{region code} so repeaters relay them only locally --- less airtime,
    same permissionless join.\\[0.4em]
    The grid-down question we opened with? \textbf{\#bitcoin-wien is already up.}
    Join it before you leave.
  \end{alertblock}
\end{frame}

\fillerslide[%
  \href{https://meshcore.io}{meshcore.io}\quad\textbullet\quad
  \href{https://meshcore-austria.at}{meshcore-austria.at}%
]{Thank you}

% ============================================================== backup / appendix
% References moved here so the main talk ends on the call to action, not a link
% dump. [noframenumbering] keeps it out of the main slide count.
\begin{frame}[noframenumbering]{References \& further reading}
  \footnotesize
  \begin{columns}[T]
    \begin{column}{0.5\textwidth}
      \textbf{Project, docs \& live}
      \begin{itemize}
        \item \href{https://meshcore.io}{meshcore.io} --- project home
        \item \href{https://docs.meshcore.io}{docs.meshcore.io} --- user \&
          protocol docs
        \item \href{https://github.com/meshcore-dev/MeshCore}{github.com/meshcore-dev/MeshCore}
          --- firmware source
        \item \href{https://logger-at.meshcore.observer}{logger-at.meshcore.observer}
          --- the Austrian packet observer
      \end{itemize}
      \medskip
      \textbf{Hardware \& radio}
      \begin{itemize}
        \item \href{https://www.semtech.com/products/wireless-rf/lora-connect/sx1262}{Semtech SX1262}
          --- the LoRa radio IC (chirp spread spectrum in silicon)
        \item \href{https://www.st.com/en/microcontrollers-microprocessors/stm32wl5x-ex.html}{ST STM32WL}
          --- the same SX126x radio IP on an MCU die
      \end{itemize}
    \end{column}
    \begin{column}{0.5\textwidth}
      \textbf{Internals \& crypto}
      \begin{itemize}
        \item \href{https://deepwiki.com/meshcore-dev/MeshCore}{deepwiki.com/meshcore-dev/MeshCore}
          --- annotated internals (incl.\ region filtering \& transport codes)
        \item \href{https://jacksbrain.com/2026/01/a-hitchhiker-s-guide-to-meshcore-cryptography/}{A
          Hitchhiker's Guide to MeshCore Cryptography} --- jacksbrain.com
        \item \href{https://github.com/michaelhart/meshcore-decoder}{michaelhart/meshcore-decoder}
          --- crypto-complete packet decoder
        \item \href{https://www.localmesh.nl/en/meshcore-encryption-details/}{localmesh.nl}
          --- encryption details (AES-128, not AES-256)
        \item \href{https://github.com/jkingsman/meshcore-hashtag-cracker}{jkingsman/meshcore-hashtag-cracker}
          --- WebGPU brute-forcer for \#hashtag-room keys
        \item \href{https://docs.meshcore.io/packet_format/}{docs.meshcore.io/packet\_format}
          --- the on-wire packet format
      \end{itemize}
    \end{column}
  \end{columns}
\end{frame}

% Optional deep-dive on the on-wire packet format --- self-contained backup
% slide after the closing, skippable on demand; ties the region code, routing
% prefix, and AES-128 + MAC together. [noframenumbering] keeps it out of the
% slide count. Verified; sources: speakers-notes.md ("Appendix 1").
\begin{frame}[noframenumbering]{Appendix 1: Anatomy of a MeshCore packet}
  \small
  \texttt{[\,header\,][\,transport codes?\,][\,path\_len\,][\,path\,][\,payload\,]}
  \hfill at most \textbf{255\,bytes} --- one LoRa frame
  \medskip
  \begin{itemize}
    \item \textbf{header} (1\,B) --- 2 bits route type, 4 bits payload type
      (advert, txt, group-text, ack, \dots), 2 bits version
    \item \textbf{transport codes} (4\,B, optional) --- only on \emph{transport}
      routes; where the 2-byte \textbf{region code} rides
    \item \textbf{path} (length byte $+$ up to 64\,B) --- \textbf{the core idea:}
      each relay \textbf{appends its own short ID} (1--3\,bytes) as it forwards
    \item \textbf{payload} (up to 184\,B) --- type-specific: an \emph{advert} is
      Ed25519-signed; a \emph{group / text} message is channel-hash $+$ 2-byte
      MAC $+$ AES-128 ciphertext
  \end{itemize}
  \vfill
  \begin{alertblock}{This is what makes it routing, not flooding}
    On a flood the path \emph{records} the relays a packet crossed; a direct
    message then \emph{carries that exact list} of router IDs, so only those
    repeaters forward it --- not every node, the way Meshtastic floods. The path
    \emph{is} the route.
  \end{alertblock}
  {\footnotesize Sources:
  \texttt{docs.meshcore.io/packet\_format}, DeepWiki 7.1, \texttt{meshcore-decoder}.}
\end{frame}

% Appendix 2: the physical layer --- what LoRa actually puts on the spectrum.
% Chirp Spread Spectrum: a packet is a preamble of up-chirps (so receivers sync),
% a sync word (down-chirps), then data symbols (chirps offset in start frequency).
% The large processing gain lets the receiver demodulate ~20 dB below the noise
% floor. Diagrams from mokosmart.com/lora-frequency (educational use, credited).
\begin{frame}[noframenumbering]{Appendix 2: What's actually on the spectrum}
  \textbf{LoRa} (the radio layer) encodes bits as \textbf{chirps} --- tones that
  sweep across the band. A packet opens with a \textbf{preamble} of up-chirps so
  receivers lock on, then a sync word, then the data symbols.
  \vfill
  \begin{columns}[b]
    \begin{column}{0.55\textwidth}\centering
      \includegraphics[width=\linewidth,height=0.40\paperheight,keepaspectratio]{lora-frequency}\\[0.3em]
      {\footnotesize Frequency vs.\ time: \textbf{preamble $\to$ sync $\to$ data}}
    \end{column}
    \begin{column}{0.43\textwidth}\centering
      \includegraphics[width=\linewidth,height=0.40\paperheight,keepaspectratio]{lora-signal-strength}\\[0.3em]
      {\footnotesize Huge \textbf{processing gain}: LoRa decodes \textbf{$\sim$20\,dB
       below the noise}. Almost magic}
    \end{column}
  \end{columns}
  \vfill
  {\scriptsize Diagrams: \href{https://www.mokosmart.com/lora-frequency/}{mokosmart.com/lora-frequency}}
\end{frame}

% Appendix 3: loop detection --- answers the question the routing slides plant
% (store-and-forward floods: why no endless echo?). Real packet storms from one
% bad repeater firmware led to loop.detect in firmware 1.14. Matrix verified
% against docs.meshcore.io/cli_commands ("Routing"); the off-by-default
% editorial is ours (re-checked against the docs 2026-08-12: still `off`).
% The flood.max.* hop budget is newer (firmware 1.16, June 2026).
% Sources & false-positive math: speakers-notes.md ("Appendix 3").
\begin{frame}[noframenumbering]{Appendix 3: Why floods don't echo forever}
  \small
  \begin{itemize}
    \item \textbf{Two brakes are built in} --- a seen-packet cache (every
      payload is hashed) means each repeater relays a flood only \emph{once},
      and the 64-byte path field caps the hops (64\,/\,32\,/\,21 at
      1\,/\,2\,/\,3-byte IDs)
    \item \textbf{One bad repeater defeats both} --- if it \emph{modifies} the
      payload before forwarding, the hash changes, the cache goes blind, and
      the same message storms the mesh as ``new'' (seen in the wild)
    \item \textbf{The v1.14 backstop:} \texttt{loop.detect} --- count
      \emph{your own ID} in the packet's path, drop past a threshold that
      scales with ID width (1-byte IDs collide --- ``me'' in the path may be
      someone else)
    \item \textbf{And a hop budget since v1.16:} \texttt{flood.max.advert}
      caps adverts at \textbf{8 hops} by default
  \end{itemize}
  \vfill
  \centering
  \renewcommand{\arraystretch}{0.92}%
  \begin{tabular}{lccc}
    \toprule
    \texttt{loop.detect} --- drop at & 1-byte IDs & 2-byte IDs & 3-byte IDs \\
    \midrule
    \texttt{off} \emph{(default)} & never & never & never \\
    \texttt{minimal}  & $\geq4\times$ & $\geq2\times$ & $\geq1\times$ \\
    \texttt{moderate} & $\geq2\times$ & $\geq1\times$ & $\geq1\times$ \\
    \texttt{strict}   & $\geq1\times$ & $\geq1\times$ & $\geq1\times$ \\
    \bottomrule
  \end{tabular}
  \vfill
  \begin{alertblock}{Off by default --- really?}
    The official reason is the transition: multibyte IDs are still rolling
    out. But \texttt{minimal} is all but false-positive-free, and one storm
    burns a whole region's duty-cycle budget --- \emph{we'd default to
    minimal}.
  \end{alertblock}
  {\footnotesize Sources: \texttt{docs.meshcore.io/cli\_commands} (Routing) ·
   blog.meshcore.io ``Path Diagnostics Improvements'' (2026-03-06).}
\end{frame}

\end{document}
